\section{Implied premium, implied risk measure and the axioms}\label{sec:axioms}

In this section we define the premium and the risk measure implied by a loss ratio and state the axioms.
Table~\ref{tab:axioms} at the end of the section summarizes them.

\subsection{The implied premium}\label{sec:implied-premium}

\begin{axiom}[\axP: one price point on each cash line]\label{ax:P}
For every $X\in L^\infty$ there is exactly one $c\in\R$ at which $\LR(X+c)\notin\R$.
\end{axiom}

\begin{definition}[Implied premium]\label{def:impliedprem}
Under \axP, the \emph{implied premium} of $X$ is $\Pi(X):=-c^*$, where $c^*$ is the unique point at which
$\LR(X+c^*)$ is not finite. We call $c^*$ the \emph{price point} of $X$.
\end{definition}

Since the cash line of $X+c$ is the cash line of $X$ shifted by $c$,
\begin{equation}\label{eq:Picash}
 \Pi(X+c)=\Pi(X)+c\qquad\text{for all }X\in L^\infty,\ c\in\R ,
\end{equation}
and $X$ belongs to the finite domain $D=\{\LR\in\R\}$ if and only if $\Pi(X)\neq0$. Axiom \axP{} does not
distinguish between the two non-finite values; they play the role of labels (Remark~\ref{rem:labels}).

The definition formalizes the notion of a price. A premium is the sure amount at which a risk changes hands:
once the insurer holds $X$ and has received $\Pi(X)$, the position $X-\Pi(X)$ commands no further premium,
and on the cash line a zero premium is exactly a point where the loss ratio is not finite. In terms of the
canonical family $\rho_\beta=(\rho-\beta\Pi)/(1-\beta)$, for books $X$ outside the break-even set,
$\rho_\beta(X)$ is the sure offset $m$ at which $\LR(X-m)=\beta$ (Corollary~\ref{cor:assoc}(b)), and $\rho_\beta\to\Pi$ as $\beta\to\infty$; the implied premium is therefore the
member of the associated family at level $\infty$, in the same way that \citet{chernymadan2009} recover a
family of risk measures from an acceptability index $\alpha$ through $\rho_x(X)=\inf\{c:\alpha(X+c)\ge x\}$.
The definition also agrees with ratemaking practice, in which the premium is obtained from a loss cost and
a permissible loss ratio \citep{wernermodlin2016}. Finally, the price point measures the premium in units of
cash, which removes the indeterminacy $(\rho,\Pi)\mapsto(\lambda\rho,\lambda\Pi)$ of the ratio. A shift
$X+c$ with $c<0$ is the book transferred together with a sure payment $|c|$.

Axiom \axP{} is ordinal: it is unchanged if $\LR$ is replaced by $\psi\circ\LR$ for any $\psi:\R\to\R$, the
non-finite values being kept. It excludes ratios that are finite along a whole cash line, such as
$\arctan(\rho/\Pi)$, and ratios with several singular points.

\subsection{Axiom \axN: a sure loss is priced at cost}

\begin{axiom}[\axN]\label{ax:N}
$\LR(c)=1$ for every constant $c\neq0$.
\end{axiom}

A sure loss carries no risk, so its fair premium is its face value; this is the ``no unjustified loading''
property of premium principles \citep{gerber1974}. Under \axP{}, axiom \axN{} implies $\Pi(c)=c$ for all $c$:
the cash line of $0$ equals $1$ away from $0$, so $\Pi(0)=0$, and $\Pi(c)=c$ by \eqref{eq:Picash}. Premiums
with a proportional loading, such as $(1+\theta)\E$, violate \axN.

\subsection{Axiom \axCM: comonotone pooling at relative cost}\label{sec:cm}

For $X,Y\in D$ with $\Pi(X)+\Pi(Y)\ne0$, where $\Pi$ is the implied premium, call
\[
 w_{X,Y}:=\frac{\Pi(X)}{\Pi(X)+\Pi(Y)},\qquad 1-w_{X,Y}=\frac{\Pi(Y)}{\Pi(X)+\Pi(Y)},
\]
the \emph{relative cost} of $X$ and of $Y$, respectively: the premium of each book as a fraction of the premiums of
the two books together.

\begin{axiom}[\axCM]\label{ax:CM}
For comonotone $X,Y\in D$ with $\Pi(X)+\Pi(Y)\neq0$,
\[
 \LR(X+Y)=w_{X,Y}\,\LR(X)+\bigl(1-w_{X,Y}\bigr)\LR(Y)
 =\frac{\Pi(X)\,\LR(X)+\Pi(Y)\,\LR(Y)}{\Pi(X)+\Pi(Y)} .
\]
\end{axiom}

Comonotone books do not hedge each other, so the cost of the pooled book should be the sum of the costs of its
parts, and then $w_{X,Y}$ is the share of $X$ in the cost of the pool. This is a consequence of the axioms and not
part of \axCM: it is the comonotonic additivity $\Pi(X+Y)=\Pi(X)+\Pi(Y)$ of Lemma~\ref{lem:comon}. Axiom \axCM{} then
says that the loss ratio of a comonotone pool is allocated to its constituents according to relative cost: the pooled
loss ratio is the combination of the two loss ratios whose weights are their relative costs. Since
$\rho=\Pi\cdot\LR$, this is the accounting statement that the loss of the pool is the sum of the losses of the
parts and the premium of the pool is the sum of their premiums, so that the loss ratio of the pool is total loss over
total premium. For books that are not comonotone, diversification will in general make the pooled premium smaller than
the sum of the premiums; this is why \citet{yaari1987} and \citet{wyp1997} impose additivity only for comonotone risks.

The weights always add up to one. They lie in $[0,1]$, and the combination is convex, exactly when $\Pi(X)$ and
$\Pi(Y)$ have the same sign, in particular when both premiums are positive, which is the case of practical interest.
Premiums of opposite sign do occur, since $\Pi(c)=c$ for sure losses and gains, and then one weight is negative and the
other exceeds one, so that the combination is affine and not convex.

A constant is comonotone with every book. Taking $Y=c$ in \axCM{} and using \axN{} and $\Pi(c)=c$ gives, for
$X\in D$, $c\ne0$ and $c\ne-\Pi(X)$,
\begin{equation}\label{eq:pool}
 \LR(X+c)=\frac{\Pi(X)}{\Pi(X)+c}\,\LR(X)+\frac{c}{\Pi(X)+c}\cdot1
 =\frac{\Pi(X)\,\LR(X)+c}{\Pi(X)+c}.
\end{equation}
Thus along the cash line the loss ratio is the combination of the loss ratio of the book and the loss ratio $1$ of the
sure loss, with the relative costs of the two as weights: adding a sure loss pulls the ratio towards $1$, and the
singularity occurs where the relative cost of the book becomes infinite, that is, at $c=-\Pi(X)$.
Proposition~\ref{prop:readings} below shows that this makes every cash line a unit-slope Möbius map.

If $\Pi(X)$ and $\Pi(Y)$ have the same sign, then $w_{X,Y}\in[0,1]$ and \axCM{} makes $\LR(X+Y)$ a convex combination of
$\LR(X)$ and $\LR(Y)$; in particular it implies the mediant inequality: $\LR(X+Y)$ lies between $\LR(X)$ and
$\LR(Y)$. The inequality is a betweenness property, and, as in decision theory, where betweenness yields only
implicit representations while weighted averages yield ratios of linear functionals
\citep{chew1983,dekel1986}, it does not suffice for our purposes (Example~\ref{ex:cube}).

Given \axP{} and \axN, \axCM{} holds if and only if the implied premium is comonotonic additive and
$\Pi\cdot\LR$, restricted to $D$, extends to a comonotonic additive functional (Lemma~\ref{lem:comon}).

\begin{remark}[A premium given from outside]\label{rem:primitive}
Suppose that a premium functional $\Pi^\circ$ is given and \axCM{} is stated with $\Pi^\circ$ as weights. If
$\LR(X)=a\ne b=\LR(Y)$ and $\LR(X+Y)=L$, then $\Pi^\circ(X)/\Pi^\circ(Y)=(b-L)/(L-a)$, so the weights are determined
by the loss ratios. On books with $\LR=1$, however, the pooled ratio equals $1$ whatever the weights, and \axCM{}
carries no information on $\Pi^\circ$. For instance, if $\LR\equiv1$ on $\{I_g\ne0\}$, then \axCM{} holds with the
non-additive weights $\Pi^\circ=(I_g)^3$, since $a^3+b^3\ne0$ implies $a+b\ne0$. With the implied premium this
difficulty does not arise, because \axP{} locates the price point also when the ratio equals $1$ elsewhere on
the cash line.
\end{remark}

\subsection{The implied risk measure}\label{sec:implied-risk}

\begin{definition}[Implied risk measure]\label{def:impliedrisk}
Under \axP, put $\rho(X):=\Pi(X)\,\LR(X)$ for $X\in D$.
\end{definition}

\begin{proposition}\label{prop:readings}
Assume \axP, \axN{} and \axCM. Then $\rho$ extends uniquely to a cash-additive functional on $L^\infty$, with
$\rho(c)=c$, and for every $X\in L^\infty$:
\begin{enumerate}[label=(\alph*),nosep]
\item $\LR(X+c)=\dfrac{\rho(X)+c}{\Pi(X)+c}$ for every $c\neq-\Pi(X)$;
\item $\rho(X)=(\Pi(X)+c)\,\LR(X+c)-c$ for every $c\neq-\Pi(X)$;
\item $\rho(X)-\Pi(X)=\lim_{c\to-\Pi(X)}(c+\Pi(X))\,\LR(X+c)$; in particular $\rho(X)=\lim_{c\to0}c\,\LR(X+c)$
  if $\Pi(X)=0$;
\item if $\rho(X)\ne\Pi(X)$, then $-\rho(X)$ is the unique zero of the cash line of $X$; if $\rho(X)=\Pi(X)$, the
  cash line equals $1$ away from the price point.
\end{enumerate}
\end{proposition}

The proof is given in Section~\ref{sec:proofs}. By (c), the implied risk measure of a book with zero premium is
the residue of its cash line at the price point: if a small premium $c$ is introduced, $\LR(X+c)\approx\rho(X)/c$.
We call $N:=\{\rho=\Pi\}$ the \emph{break-even set}; $N\cap D=\{\LR=1\}$, and $N$ also contains the books with
$\rho=\Pi=0$. On $N$ the residue vanishes and, by (d), the cash line has no zero, so that the risk measure cannot
be read from the zero of the cash line; this is why it is defined through (b).

\begin{example}[Three cash lines]\label{ex:three}
Let $\rho=I_{\mathrm{id}}=\E$ and $\Pi=I_{\sqrt{\cdot}}$, and let $P(A)=\tfrac14$. Figure~\ref{fig:three} shows
the cash lines of three books.
\begin{enumerate}[label=(\roman*),nosep]
\item $X=\1_A$: $\LR(X+c)=(\tfrac14+c)/(\tfrac12+c)$. The price point is $-\tfrac12$, so $\Pi(X)=\tfrac12$, and
  $\rho(X)=\Pi(X)\LR(X)=\tfrac14$; the cash line vanishes at $-\tfrac14$.
\item $Y=\1_A-\tfrac12$, a book with zero premium: $\LR(Y)=\infty$ and $\LR(Y+c)=(c-\tfrac14)/c$. The residue at
  the price point $0$ is $-\tfrac14=\rho(Y)$.
\item $Z\equiv\tfrac12$, a break-even book: $\LR(Z+c)=1$ for $c\ne-\tfrac12$, and $\LR(Z-\tfrac12)=\star$. The
  price point gives $\Pi(Z)=\tfrac12$; the cash line has no zero, and $\rho(Z)=\Pi(Z)\cdot1=\tfrac12$.
\end{enumerate}
\end{example}

\begin{figure}[ht]
\centering
\begin{tikzpicture}
\begin{axis}[width=0.8\textwidth,height=6.4cm,xmin=-1.2,xmax=1.2,ymin=-4,ymax=4,
  axis lines=middle,xtick={-1,-0.5,0,0.5,1},xlabel={$c$},ylabel={$\LR(\cdot+c)$},
  xlabel style={at={(axis cs:1.2,0)},anchor=west},ylabel style={at={(axis cs:0,4)},anchor=south},
  samples=150,unbounded coords=jump,tick label style={font=\small},
  legend style={font=\small,at={(0.98,0.03)},anchor=south east,cells={anchor=west}}]
\addplot[thick,blue!70!black,domain=-1.2:-0.53] {(0.25+x)/(0.5+x)};
\addlegendentry{(i) $X=\1_A$}
\addplot[thick,red!70!black,dashed,domain=-1.2:-0.04] {(x-0.25)/x};
\addlegendentry{(ii) $Y=\1_A-\frac12$}
\addplot[thick,green!45!black,dotted,domain=-1.2:-0.52] {1};
\addlegendentry{(iii) $Z\equiv\frac12$}
\addplot[thick,blue!70!black,domain=-0.47:1.2,forget plot] {(0.25+x)/(0.5+x)};
\addplot[thick,red!70!black,dashed,domain=0.04:1.2,forget plot] {(x-0.25)/x};
\addplot[thick,green!45!black,dotted,domain=-0.48:1.2,forget plot] {1};
\draw[gray!70] (axis cs:-0.5,-4) -- (axis cs:-0.5,4);
\addplot[only marks,mark=*,mark size=1.8pt,blue!70!black,forget plot] coordinates {(-0.25,0)};
\addplot[only marks,mark=*,mark size=1.8pt,red!70!black,forget plot] coordinates {(0.25,0)};
\addplot[only marks,mark=o,mark size=2pt,green!45!black,forget plot] coordinates {(-0.5,1)};
\end{axis}
\end{tikzpicture}
\caption{Cash lines of the three books of Example~\ref{ex:three}. Book (i) explodes at $c=-\frac12$ (its price
point) and vanishes at $c=-\frac14$; book (ii) explodes at $c=0$ with residue $-\frac14$; book (iii) equals $1$
except at its price point $c=-\frac12$, where the singularity is removable (open circle).}\label{fig:three}
\end{figure}

\begin{remark}[Labels at the price point]\label{rem:labels}
To have $\LR=\bracket{\rho}{\Pi}$ also at price points, one may adopt the convention
\begin{center}
\ax{S}\quad at the price point, $\LR=\star$ if the residue in (c) is zero and $\LR=\infty$ otherwise.
\end{center}
Convention \ax{S} is not used in the proofs. If, instead, the loss ratio is real valued with $0/0:=1$, the price
point of a break-even book cannot be observed and the implied premium is not defined there.
\end{remark}

\begin{remark}[A formulation by cash lines]\label{rem:oldCP}
One may alternatively postulate (CP): every cash line is a unit-slope Möbius map,
$\LR(X+c)=\bracket{\rho(X)+c}{\Pi(X)+c}$. Proposition~\ref{prop:readings}(a) shows that (CP) follows from \axP,
\axN{} and \axCM, the labels agreeing under \ax{S}. Conversely, (CP) implies \axP{} and \ax{S}, and (CP) together
with \axCM{} implies \axN: if $\pi_0:=\Pi(0)\ne0$, \axCM{} applied to the constants $a=b=-\pi_0/2$ forces
$\LR(-\pi_0)\in\R$, although $-\pi_0$ is the price point of $0$; hence $\pi_0=0$, and \axCM{} with $X=Y=a\ne0$ gives
$(\rho(0)+2a)/(2a)=(\rho(0)+a)/a$, i.e.\ $\rho(0)=0$ and $\LR(c)=1$. Thus \axP+\axN+\axCM+\ax{S} is equivalent to
(CP)+\axCM.
\end{remark}

\subsection{Axiom \axLI: law invariance}\label{sec:li}

\begin{axiom}[\axLI]\label{ax:LI}
If $X\eqd Y$, then either $\LR(X)=\LR(Y)$ or both values are non-finite.
\end{axiom}

This is the standard assumption of premium and risk-measure theory \citep{wyp1997,kusuoka2001}. Because the
implied premium and risk measure are computed from cash lines, \axLI{} passes to them (Lemma~\ref{lem:LI}).
Ratios that weight scenarios non-uniformly, such as $\E_Q/\E$ with $Q\neq P$, violate \axLI; they are covered by
Theorem~\ref{thm:A0}.

\subsection{Axioms \axRC{} and \axRM: regularity of the implied pair}\label{sec:reg}

Under \axP, write $\rho_D:=\Pi\cdot\LR$ on $D$. The regularity axioms are stated for $\Pi$ and $\rho_D$, so that
they make sense independently of the other axioms.

\begin{axiom}[\axRC]\label{ax:RC}
$\Pi$ is uniformly continuous on $(L^\infty,\|\cdot\|_\infty)$ and $\rho_D$ is uniformly continuous on $D$.
\end{axiom}

\begin{axiom}[\axRM]\label{ax:RM}
If $X\le Y$ almost surely, then $\Pi(X)\le\Pi(Y)$, and $\rho_D(X)\le\rho_D(Y)$ if $X,Y\in D$.
\end{axiom}

Under \axN{} and \axCM{} these axioms are equivalent to the corresponding properties of the implied risk
measure on $L^\infty$, and \axRM{} implies \axRC{} (Lemma~\ref{lem:cashlip}). We include the signed tier
(\axRC) because the canonical family $\rho_\beta$ is a difference of monotone functionals and is in general not
monotone (Corollary~\ref{cor:assoc} and Example~\ref{ex:family}). Monotonicity is imposed on the implied pair and
not on the ratio, for the following reason.

\begin{proposition}[No monotone loss ratio]\label{prop:F2}
Let $h,g\in\Hset_1$ with $g>0$ on $(0,1]$, and let $\LR=I_h/I_g$ on $\Xp:=L^\infty_+\setminus\{0\}$.
If $\LR$ is monotone on $\Xp$ (non-decreasing or non-increasing for the almost sure order),
then $h=g$ and $\LR\equiv1$.
\end{proposition}

\begin{proof}
Let $P(A)=p>0$. For $c\ge0$, $\LR(\1_A+c)=\frac{h(p)+c}{g(p)+c}$ is monotone in $c$ with derivative at
$0$ of the sign of $g(p)-h(p)$. Next let $P(E)=q>0$ and $t\ge0$. Then $X_t=1+t\1_E\in\Xp$ is
increasing in $t$, and $I_k(X_t)=1+t\,k(q)$ for $k\in\{h,g\}$, so
$\LR(X_t)=\frac{1+t\,h(q)}{1+t\,g(q)}$ has derivative at $t=0$ of the sign of $h(q)-g(q)$. If $\LR$
is non-decreasing both signs are non-negative, hence $h=g$ on $(0,1]$; if $\LR$ is non-increasing both are
non-positive, and again $h=g$. \qedhere
\end{proof}

\begin{example}\label{ex:F2}
Let $h=\mathrm{id}$, $g=\sqrt{\cdot}$ and $p=q=\tfrac14$ in the proof. Both families of books in
Figure~\ref{fig:F2} increase in the almost sure order, but the loss ratio increases along the first and
decreases along the second: adding a sure loss to an over-priced book moves the ratio towards $1$ from below,
whereas adding risk to a sure loss moves it away from $1$.
\end{example}

\begin{figure}[ht]
\centering
\begin{tikzpicture}
\begin{axis}[width=0.7\textwidth,height=5.4cm,xmin=0,xmax=8,ymin=0.4,ymax=1.05,
  xlabel={$s$},ylabel={loss ratio},tick label style={font=\small},samples=100,domain=0:8,
  legend style={font=\small,at={(0.98,0.5)},anchor=east,cells={anchor=west}}]
\addplot[thick,blue!70!black] {(0.25+x)/(0.5+x)};
\addlegendentry{$\LR(\1_A+s)$}
\addplot[thick,red!70!black,dashed] {(1+0.25*x)/(1+0.5*x)};
\addlegendentry{$\LR(1+s\1_E)$}
\addplot[gray,dotted,domain=0:8] {1};
\end{axis}
\end{tikzpicture}
\caption{Example~\ref{ex:F2}: $h=\mathrm{id}$, $g=\sqrt{\cdot}$, $P(A)=P(E)=\frac14$. Both families of books
increase with $s$; the loss ratio increases along one and decreases along the other.}\label{fig:F2}
\end{figure}

Consequently, the theorems on acceptability indices \citep{chernymadan2009}, which require monotonicity of the
index, do not apply to loss ratios.

\subsection{Summary}

\begin{table}[ht]
\centering\small
\begin{tabularx}{\textwidth}{@{}l>{\raggedright\arraybackslash}p{4.2cm}>{\raggedright\arraybackslash}X@{}}
\toprule
\textbf{Axiom} & \textbf{Statement} & \textbf{Interpretation and counterpart}\\
\midrule
\axP & one price point on each cash line & defines the premium; the level $\beta=\infty$ of the associated family\\
\axN & $\LR(c)=1$ for $c\ne0$ & no unjustified loading \citep{gerber1974}\\
\axCM & comonotone pooling at relative cost & comonotonic additivity \citep{yaari1987,wyp1997}\\
\axLI & law invariance & standard \citep{wyp1997,kusuoka2001}\\
\axRC/\axRM & regularity of the implied pair & continuity / monotonicity of premia\\
\bottomrule
\end{tabularx}
\caption{The axioms.}\label{tab:axioms}
\end{table}
